%% kapa: exam=1 round=34 year=2023 period=1 form=A subject=civil_law subject_ko=민법 seq=2 number=2
%% source: Q-Net 감정평가사 기출문제 (cst00309) artlSeq=5211789 file=raw/1008/5211789/2023년도 제34회 감정평가사 1차시험 1교시.pdf
%% conversion: kordoc 4.24.51 -> tools/gen1.py (auto); review=proofread-sonnet; flags=none
\begin{question}{2}
제한능력에 관한 설명으로 옳지 않은 것은? (다툼이 있으면 판례에 따름)
\begin{choices}
  \item 성년후견인은 여러 명을 둘 수 있다.
  \item 가정법원은 본인의 청구에 의하여 취소할 수 없는 피성년후견인의 법률행위의 범위를 변경할 수 있다.
  \item 가정법원이 피성년후견인에 대하여 한정후견개시의 심판을 할 때에는 종전의 성년후견의 종료 심판을 하여야 한다.
  \item 한정후견의 개시를 청구한 사건에서 의사의 감정 결과 성년후견 개시의 요건을 충족하고 있다면 법원은 본인의 의사를 고려하지 않고 성년후견을 개시할 수 있다.
  \item 특정후견의 심판이 있은 후에 피특정후견인이 특정후견인의 동의 없이 재산상의 법률행위를 하더라도 이는 취소의 대상이 되지 않는다.
\end{choices}
\end{question}
\answer{4}
